\documentclass[../../main.tex]{subfiles}
\begin{document}

\begin{flushright}
\footnotesize\textit{【自動文字起こし・要確認】}
\end{flushright}

{\bf [解]}
\paragraph{(1)}

合同式の法を$p$とすると，$a_n\equiv b_n$である．漸化式から
\begin{align}
b_{n+2}\equiv b_{n+1}(b_n+1)\quad\cdots①
\end{align}
だから，題意は示された．$\blacksquare$

\paragraph{(2)}

$b_1=2$，$b_2=3$であるから，①より，
\begin{align}
b_3\equiv3\cdot3=9,\quad b_4\equiv9\cdot4=36\equiv2,\quad b_5\equiv2\cdot10\equiv3\pmod{17}
\end{align}
だから，以降$\{2,3,9\}$のくり返しとなるから，
\begin{align}
b_1=b_4=b_7=b_{10}=2,\qquad b_2=b_5=b_8=3,\qquad b_3=b_6=b_9=9
\end{align}

\paragraph{(3)}

背理法で示す．$b_n\ne b_m$と仮定する．この時，$0\le b_n<b_m\le p-1$としてよい．条件から
\begin{align}
b_{n+2}&\equiv b_{n+1}(b_n+1)\\
b_{m+2}&\equiv b_{m+1}(b_m+1)
\end{align}
が成り立つ．$r=b_m-b_n$とおくと，$r\in\mathbb N$かつ$1\le r\le p-1$であり，$b_{n+1}=b_{m+1}$，$b_{n+2}=b_{m+2}$を用いて上の2式を辺々引いて，
\begin{align}
0\equiv b_{n+1}\cdot r\pmod p
\end{align}
となる．しかし，$1\le b_{n+1}\le p-1$，$1\le r\le p-1$であり$p$は素数だから，これは矛盾．よって$b_n=b_m$である．$\blacksquare$

\paragraph{(4)}

題意が成立する時，任意の$n\in\mathbb N_{\ge2}$に対して，$b_n\ne0$が成立する．さらに，この時，$b_n$（$n\ge2$）は$1,2,\ldots,p-1$の高々$p-1$個の値しかとらないことから，連続する2項$b_n,b_{n+1}$のならびは高々$(p-1)^2$通りしか存在せず，
\begin{align}
b_{n+2}=b_{m+2}>0,\quad b_{n+1}=b_{m+1}>0
\end{align}
を満たす異なる$n,m\in\mathbb N_{\ge2}$がある（$n<m$とする）．この時(3)から$b_n=b_m$となり，再び$b_{n+1}=b_{m+1}>0$，$b_n=b_m>0$に(3)を適用して，$b_{n-1}=b_{m-1}$となる．これをくり返して，$n<m$より，
\begin{align}
b_1=b_{m+1-n}>0\quad(\because b_n>0)
\end{align}
となるから，$b_1\ne0$，つまり$a_1$は$p$で割り切れない．$\blacksquare$

\bigskip
\noindent{\bf [解2]}
\paragraph{(4)}

[解]と同様に，
\begin{align}
b_t=b_s,\qquad b_{t+1}=b_{s+1}\qquad(s<t,\ s,t\in\mathbb N_{\ge2})
\end{align}
を満たす$s,t$がある．このうち，$s$が最小であるものについて考えると，$a_2,a_3,\ldots$が$p$で割り切れないことから，$b_s\ne0$，したがって$b_t=b_s\ne0$である．$s>1$と仮定すると，$b_{s-1},b_{t-1}$が定義され，$b_s=b_t>0$かつ$b_{s+1}=b_{t+1}$だから，(3)より$b_{s-1}=b_{t-1}$となり，$s-1<t-1$は$s$の最小性に反する．したがって$s=1$で，
\begin{align}
b_1=b_s=b_t\ne0
\end{align}
（$\because t\ge2$かつ$a_2,a_3,\ldots$が$p$で割り切れない）だから，$a_1$は$p$で割り切れない．$\blacksquare$

\end{document}
